% TOP_PROBLEM: {"id": 3202, "title": "Edge list coloring conjecture", "category_id": 3, "category": {"id": 3, "name": "graph_theory", "display_name": "Graph Theory", "description": "Problems involving graphs, networks, and their properties.", "slug": "graph-theory", "order_index": 3, "created_at": "2026-07-31T15:26:25.671Z"}, "status": "open", "problem_number": "OPG-2110", "set_id": 12}
% FIRST_POSTED: 2026-10-01
% ENTRY_KIND: statement_only
% UnsolvedMath ID 3202; source catalogue code OPG-2110
% !TeX program = lualatex
% Definition and sources only; this file is not an LLM solution attempt.
\documentclass[11pt,a4paper]{article}
\usepackage[margin=25mm]{geometry}
\usepackage{fontspec}
\setmainfont{lmroman10-regular.otf}[BoldFont=lmroman10-bold.otf,
 ItalicFont=lmroman10-italic.otf,BoldItalicFont=lmroman10-bolditalic.otf]
\usepackage{amsmath,amssymb,mathtools}
\usepackage{unicode-math}
\setmathfont{latinmodern-math.otf}
\AtBeginDocument{\let\setminus\smallsetminus}
\usepackage{xurl,hyperref}
\hypersetup{colorlinks=true,urlcolor=blue}
\setcounter{secnumdepth}{0}
\setlength{\parindent}{0pt}
\setlength{\parskip}{5pt}
\setlength{\emergencystretch}{3em}
\begin{document}
\section*{Edge list coloring conjecture}\label{3202}
{\small UnsolvedMath ID 3202 · OPG-2110 · Graph Theory · OpenGarden}
\subsection{Definitions and mathematical statement}
Let $G=(V,E)$ be a finite undirected loopless
multigraph. Parallel edges are distinguished
and share both endpoints. A proper edge
colouring assigns a colour to each edge
so that different edges sharing an endpoint
receive different colours. The edge chromatic
number, or chromatic index, $\chi'(G)$ is
the minimum number of colours in such a
colouring.

An edge-list assignment specifies a finite
set $L(e)$ of allowed colours for every
$e\in E$. An $L$-edge-colouring is a proper
edge colouring $c$ with $c(e)\in L(e)$
for all $e$. Define $\chi'_\ell(G)$ to
be the least nonnegative integer $q$ such
that every assignment with $|L(e)|\geq q$
for all edges has an $L$-edge-colouring.
For an edgeless graph both parameters
are zero.

The conjecture is
\[
 \chi'_\ell(G)=\chi'(G)
 \qquad\text{for every finite loopless multigraph }G.
\]
The lists are arbitrary: they need not
be equal, and their union can be much
larger than $\chi'(G)$. Assigning the same
$q$-element list to every edge immediately
gives $\chi'_\ell(G)\geq\chi'(G)$; the
opposite inequality is the requested result.

Equivalently, ordinary chromatic number and
list chromatic number should coincide on
line graphs of loopless multigraphs, whose
vertices represent edges and whose adjacency
records common endpoints. The conjecture
does not request a bound solely in terms
of maximum degree.
\subsection{Short English statement}
Are as many colours per arbitrary edge list as the ordinary chromatic index always sufficient for proper edge colouring?
\subsection{Sources}
\begin{itemize}
\item[S1] ULAM.AI and the UnsolvedMath Contributors, supplied problems.json, record 3202 (OPG-2110), OpenGarden collection. Catalogue identity and source transcription; CC BY 4.0.
\url{https://huggingface.co/datasets/ulamai/UnsolvedMath}.
\item[S2] Open Problem Garden, Edge list coloring conjecture. Original problem and discussion; the mathematical formulation below is expanded from the supplied source record.
\url{http://www.openproblemgarden.org/op/edge_list_coloring_conjecture}.
\end{itemize}
\end{document}
